\section{Generative Adversarial Networks (GANs)}

\subsection{从显式建模到隐式采样}

VAE 试图\textbf{显式地}建模数据的概率密度 $p(x)$。但如果我们的目标只是\textbf{生成新样本}，是否一定需要显式建模密度？

\begin{figure}[H]
\centering
\includegraphics[width=0.85\textwidth]{slides-images/slide-48.jpg}
\caption{GAN 的思路：不显式建模密度，直接学习从噪声到数据的映射\protect\footnotemark}
\end{figure}
\footnotetext{Slides 第48页。}

\begin{knowledgebox}{GAN 的设计哲学}
GAN 的核心思路是：不去显式建模 $p(x)$，而是直接学习一个从简单分布（如高斯噪声）到复杂数据分布的\textbf{变换}。具体做法是从简单的噪声分布中采样，然后训练一个生成器网络将噪声变换为逼真的数据样本。关键问题变成了：如何训练这个生成器？
\end{knowledgebox}

\subsection{GAN 的对抗训练思想}

GAN 提出了一个优雅的解决方案：引入两个互相竞争的网络——\textbf{生成器}（Generator, $G$）和\textbf{判别器}（Discriminator, $D$）。

\begin{itemize}
    \item \textbf{生成器 $G$}：接收随机噪声 $z$，尝试生成尽可能逼真的假数据 $x_{\text{fake}} = G(z)$
    \item \textbf{判别器 $D$}：接收数据样本（可能是真实的或生成的），输出该样本为真实数据的概率 $D(x) \in {[}0, 1{]}$
\end{itemize}

\begin{figure}[H]
\centering
\includegraphics[width=0.85\textwidth]{slides-images/slide-68.jpg}
\caption{GAN 的训练架构：生成器与判别器的对抗\protect\footnotemark}
\end{figure}
\footnotetext{Slides 第68页。}

两个网络形成\textbf{对抗博弈}：
\begin{itemize}
    \item $G$ 试图生成足以欺骗 $D$ 的假数据
    \item $D$ 试图正确区分真实数据和 $G$ 生成的假数据
    \item 在训练达到全局最优时，$G$ 能够生成与真实数据分布一致的样本
\end{itemize}

\subsection{GAN 训练的直觉}

课程用一维数据线上的动画生动展示了 GAN 的训练过程：

\begin{figure}[H]
\centering
\includegraphics[width=0.85\textwidth]{slides-images/slide-50.jpg}
\caption{GAN 训练直觉：生成器从噪声出发，逐步逼近真实数据分布\protect\footnotemark}
\end{figure}
\footnotetext{Slides 第50页。}

\textbf{训练过程的迭代}：
\begin{enumerate}
    \item 生成器初始时从随机噪声生成质量很差的假数据
    \item 判别器学习区分真实数据（绿色点）和假数据（粉色点），为真实数据赋予高概率
    \item 生成器根据判别器的反馈，调整生成的数据分布，使假数据更接近真实数据
    \item 判别器再次更新以适应更好的假数据
    \item 反复迭代，直到假数据与真实数据无法区分
\end{enumerate}

\begin{figure}[H]
\centering
\includegraphics[width=0.85\textwidth]{slides-images/slide-65.jpg}
\caption{训练后期：假数据（粉色）与真实数据（绿色）分布接近\protect\footnotemark}
\end{figure}
\footnotetext{Slides 第65页。}

\subsection{GAN 的损失函数}

GAN 的训练目标可以用一个 \textbf{minimax} 博弈形式表达：

$$\min_G \max_D \; \mathbb{E}_{z,x}\left[\log D(x) + \log(1 - D(G(z)))\right]$$

\begin{figure}[H]
\centering
\includegraphics[width=0.85\textwidth]{slides-images/slide-70.jpg}
\caption{GAN 损失函数：生成器最小化、判别器最大化同一目标\protect\footnotemark}
\end{figure}
\footnotetext{Slides 第70页。}

\begin{importantbox}{GAN 损失函数解读}
\textbf{判别器的目标（最大化）}：
\begin{itemize}
    \item $\log D(x)$：对真实数据 $x$，希望 $D(x) \to 1$（正确识别为真）
    \item $\log(1 - D(G(z)))$：对生成数据 $G(z)$，希望 $D(G(z)) \to 0$（正确识别为假）
\end{itemize}

\textbf{生成器的目标（最小化）}：
\begin{itemize}
    \item 希望 $D(G(z)) \to 1$，即让判别器将假数据误判为真
    \item 等价于最小化 $\log(1 - D(G(z)))$
\end{itemize}
\end{importantbox}

\begin{warningbox}{GAN 训练的不稳定性}
GAN 的对抗训练在实践中可能非常不稳定。两个网络需要保持相对平衡的训练进度——如果判别器太强，生成器得不到有用的梯度信号；如果生成器太强，判别器无法提供有效的反馈。常见问题包括模式崩塌（Mode Collapse）和训练振荡。后续工作如 Wasserstein GAN (WGAN) 提出了更稳定的替代损失函数。
\end{warningbox}

\subsection{GAN 作为分布变换器}

训练完成后，我们丢弃判别器，只使用生成器来生成新数据：

\begin{figure}[H]
\centering
\includegraphics[width=0.85\textwidth]{slides-images/slide-72.jpg}
\caption{训练后的 GAN：仅使用生成器从噪声创建新数据\protect\footnotemark}
\end{figure}
\footnotetext{Slides 第72页。}

从根本上讲，GAN 的生成器学习了一个\textbf{分布变换}：从高斯噪声分布 $z \sim \mathcal{N}(0, 1)$ 到目标数据分布的映射。

\begin{figure}[H]
\centering
\includegraphics[width=0.85\textwidth]{slides-images/slide-75.jpg}
\caption{GAN 是分布变换器：从噪声分布到数据分布的映射\protect\footnotemark}
\end{figure}
\footnotetext{Slides 第75页。来源：Brock+ ICLR 2019。}

在噪声空间中进行平滑插值时，生成的样本也会在数据流形上平滑过渡，说明生成器学到了数据的连续结构。

\subsection{GAN 的发展与应用}

\textbf{Progressive Growing of GANs}：通过逐步增加网络层数和分辨率，可以生成极为逼真的高分辨率人脸图像。

\begin{figure}[H]
\centering
\includegraphics[width=0.85\textwidth]{slides-images/slide-78.jpg}
\caption{Progressive GAN 生成的高分辨率人脸（均为 AI 生成）\protect\footnotemark}
\end{figure}
\footnotetext{Slides 第78页。来源：Karras+ ICLR 2018。}

\textbf{Conditional GAN 与 Pix2Pix}：通过在生成器和判别器中加入条件信息，可以实现配对的图像翻译任务——如从语义分割图生成真实街景照片。

\begin{figure}[H]
\centering
\includegraphics[width=0.85\textwidth]{slides-images/slide-80.jpg}
\caption{Conditional GAN (Pix2Pix)：配对图像翻译\protect\footnotemark}
\end{figure}
\footnotetext{Slides 第80页。来源：Isola+ CVPR 2017。}

\begin{figure}[H]
\centering
\includegraphics[width=0.85\textwidth]{slides-images/slide-82.jpg}
\caption{Pix2Pix 应用结果：地图与航拍图的双向转换\protect\footnotemark}
\end{figure}
\footnotetext{Slides 第82页。来源：Isola+ CVPR 2017。}

\textbf{CycleGAN}：不需要配对数据即可实现跨域转换。不仅限于图像，还可以应用于语音变换等领域。

\begin{figure}[H]
\centering
\includegraphics[width=0.85\textwidth]{slides-images/slide-85.jpg}
\caption{CycleGAN 在语音变换中的应用\protect\footnotemark}
\end{figure}
\footnotetext{Slides 第85页。}

\begin{knowledgebox}{CycleGAN 的创新}
传统 GAN 从噪声分布生成目标分布。CycleGAN 扩展了这个思想：它不从噪声出发，而是从一个数据分布出发，学习到另一个数据分布的映射。关键创新是使用\textbf{循环一致性损失}（Cycle Consistency Loss），确保从 A 域转到 B 域再转回 A 域时，能恢复原始输入。这使得无需配对数据即可训练。
\end{knowledgebox}

\subsection{本章小结}

GAN 通过生成器-判别器的对抗训练，实现了不需要显式建模密度即可生成高质量样本的目标。GAN 本质上是一个分布变换器，将简单的噪声分布映射到复杂的数据分布。虽然 GAN 训练存在不稳定的挑战，但其生成质量（尤其是图像）通常优于 VAE，并且衍生出了 Conditional GAN、CycleGAN 等强大的变体。


